\clearpage\NoSourceHead\SourceLeaf{407}
\SourceGroup{V}
\HistoricTitle{CONVERGENZ DER REIHEN, IN WELCHE DIE PERIODISCHEN FUNCTIONEN EINER VERÄNDERLICHEN GRÖSSE ENTWICKELT WERDEN.}
\HistoricCitation{[Handschrift von {\scriptsize 13} mit A{\scriptsize 1} bis A{\scriptsize 13} bezeichneten Blättern in Fa, Kapsel {\scriptsize 46}a.]}
\noindent\hspace*{\parindent}Ehe ich zu der Untersuchung selbst, welcher diese Denkschrift gewidmet ist, übergehe, verweile ich etwas bei einigen Begriffsbestimmungen, die anderwerts entweder noch gar nicht vorkommen, oder nicht überall auf die gleiche Weise angewandt werden.
\begin{center}1.\par\smallskip\bfseries Convergenz der unendlichen Reihen im Allgemeinen.\end{center}
\noindent\hspace*{\parindent}Ich werde unter \textls[120]{Convergenz}, einer unendlichen Reihe schlechthin beigelegt, nichts anderes verstehen, als die beim unendlichen Fortschreiten der Reihe eintretende unendliche Annäherung ihrer Glieder an die 0. Die Convergenz einer Reihe an sich ist also wohl zu unterscheiden von der Convergenz ihrer Summirung zu einem endlichen Grenzwerthe; letztere schliesst zwar die erstere ein, aber nicht umgekehrt.

Divergent heisst eine unendliche Reihe, in welcher bei hinlänglich weiter Fortsetzung Glieder erscheinen, die jede (vorher) gegebene Grösse überschreiten.

Reihen, die weder schlechthin convergent, noch divergent sind, werden entweder in ihrem unendlichen Fortschreiten eine unendliche Annäherung an einen von 0 verschiedenen Grenzwerth darbieten, in welchem Fall man sie convergent zu diesem Werthe nennen kann, oder so beschaffen sein, dass sie
\clearpage\SourceLeaf{408}\ArithmeticHead{401}{R}{CONVERGENZ DER ENTWICKLUNG PERIODISCHER FUNCTIONEN.}
\noindent nie aufhören in einem endlichen Zwischenraum hin und her zu schwanken. Für unsern Zweck ist es jedoch nicht nothwendig, diese Unterscheidung noch weiter ins Einzelne zu verfolgen.

Wir haben demnach drei Arten oder Ordnungen von unendlichen Reihen, oder da zwei dieser Arten wieder in je zwei zerfallen, fünf Ordnungen, nemlich:

I. Divergente Reihen.

II. Zwischen gewissen Grenzen stets schwankend bleibende Reihen.

Endlich die zu einem bestimmten Grenzwerthe convergirende[n] Reihen, nemlich

III. Reihen, die zu einem von 0 verschiedenen Grenzwerthe convergiren.

IV. und V. Die zu 0 convergirenden Reihen, welche ich ohne weitern Zusatz convergirende Reihen nenne. In die V. Classe setze ich diejenigen, bei denen die successive Summation eine unendliche Annäherung an eine endliche bestimmte Grenze darbietet, in die IV. die, wo dies nicht Statt findet.

Durch Addition zweier Reihen zu einander (indem man das erste Glied der einen Reihe zum ersten der andern addirt, das zweite zum zweiten u.s.w.) entsteht eine neue Reihe, über deren Classe sich folgendes sagen lässt. Die Addition zweier Reihen der fünften Classe gibt eine der fünften; die Addition einer der fünften zu einer der vierten, gibt eine der vierten, und allgemein bleibt die Ordnung ungeändert dieselbe, wenn eine Reihe der fünften oder überhaupt eine von höherer Ordnung hinzuaddiert wird. Man schliesst daraus leicht allgemein, dass durch Addition beliebig vieler Reihen eine Reihe entsteht, deren Ordnung mit der niedrigsten Ordnung übereinstimmt, wenigstens allemahl, wenn von dieser niedrigsten Ordnung nur eine unter den addierten ist: die Addition zweier oder mehrerer von gleicher Ordnung kann allerdings in speciellen Fällen eine von höherer hervorbringen. Hiemit hängt auch die Classification derjenigen Reihen zusammen, deren Glieder complexe Grössen sind. Man betrachtet die Reihe der reellen Theile, und die der imaginären jede für sich: gehören sie zu gleicher Ordnung, so setzt man auch die ganze Reihe in eben dieselbe, im entgegengesetzten Falle richtet sich letztere nach der niedrigern Stufe.

Wenn man die Glieder einer unendlichen Reihe
\[
R=A,A',A'',A'''\ \text{u.s.w.}
\]
\par\vfill\noindent{\small X1.}\hfill{\small 51}\par
\clearpage\SourceLeaf{409}\ArithmeticHead{402}{L}{ANALYSIS. NACHLASS.}
\noindent mit den Gliedern einer geometrischen Reihe
\[
1,h,hh,h^3\ \text{u.s.w.}
\]
\noindent multiplicirt, so wird die neue Reihe
\[
S=A,hA',hhA'',h^3A'''\ \text{u.s.w.}
\]
\noindent in dem Falle, wo \(h\) kleiner ist als 1, nicht bloss convergent sein, sondern auch eine convergente Summation darbieten, also zur 5\,ten Ordnung gehören, allemahl wo \(R\) nicht divergent ist, also zur 2., 3., 4. oder 5. Ordnung gehört. In der That, da in jeder nicht divergierenden Reihe keines der Glieder eine gewisse Grösse \(g\)[\textsuperscript{*})] überschreiten kann, so wird in der Reihe \(S\) das Glied mit dem Index \(n\) jedenfalls nicht grösser sein als \(gh^n\), und jedes der folgenden noch kleiner, und eben so wird die Summation der auf dies Glied folgenden, wie weit man sie auch fortsetzt, nicht grösser sein als
\[
\frac{gh^n}{1-h},
\]
\noindent welche Grösse einen so kleinen nicht verschwindenden Werth \(c\) als man nur wünscht, erlangen kann, wenn man nur \(n\) gross genug annimmt.

Ist \(h\) grösser als 1, oder werden die Glieder der Reihe \(R\) mit den Gliedern einer steigenden geometrischen Progression multiplicirt, so sind drei Fälle denkbar:

I. Indem man \(h\) als veränderlich betrachtet und vom Werthe [Eins] an allmählig wachsen lässt, hört \(S\) sofort auf convergent zu sein, wie wenig auch \(h\) den Werth 1 überschreitet.

II. \(S\) bleibt convergent bis zu einem bestimmten Werte von \(h\), nemlich \(h=\frac1e\), wo \(e\) ein ächter Bruch ist. Dies bis kann einschliesslich oder ausschliesslich zu verstehn sein, nemlich entweder ist \(\frac1e\) der letzte Werth des wachsenden \(h\), für welchen \(S\) convergent ist, oder der erste für welchen \(S\) nicht convergent ist.

III. \(S\) bleibt stets convergent, wie gross man auch \(h\) annehmen möge, wie z. B. in dem Falle, wo \(R\) eine Exponentialgrösse oder den Sinus oder Cosinus eines Bogens ausdrückt.
\par\vfill\noindent\rule{.25\linewidth}{.4pt}\par\smallskip
{\footnotesize\noindent\hspace*{\parindent}[\textsuperscript{*}) In der Handschrift steht an dieser Stelle noch das Wort »nicht«.]\par}
\clearpage\SourceLeaf{410}\ArithmeticHead{403}{R}{CONVERGENZ DER ENTWICKLUNG PERIODISCHER FUNCTIONEN.}
\noindent\hspace*{\parindent}Im Falle I. convergirt also die Reihe \(R\) langsamer, im Falle III. schneller als jede fallende geometrische Progression; so wie im Falle II. jene Reihe langsamer convergirt als jede fallende geometrische Progression mit einem kleinern Exponenten als \(e\), aber schneller als jede mit irgend einem grössern Exponenten als \(e\). In so fern steht also unter allen fallenden geometrischen Progressionen, diejenige deren Exponent \(=e\) ist, in Beziehung auf Convergenz der Reihe \(R\) \textls[120]{am nächsten}, und es wird daher nicht unschicklich sein, \(e\) den Exponenten der Convergenz der Reihe \(R\) zu nennen. Bei alledem ist klar, dass je nachdem \(S\) selbst, für \(h=\frac1e\), noch convergent oder schon divergent wird, jene Progression noch weniger oder mehr convergent ist als \(R\), und wir werden weiter unten zeigen, wie eine noch nähere Anschmiegung erreicht werden kann[\textsuperscript{*})].
\begingroup\makeatletter
\long\def\@makefntext#1{\noindent\hspace*{\parindent}#1}
\footnotetext{[\textsuperscript{*}) Siehe den unten folgenden \textls[120]{Anhang}.]}
\endgroup
